\chapter{Mathematical Preliminaries}\label{c1}
These notes follow the chapters of \cite{so1989} as a general route towards 
understanding methods of quantum chemistry. The book was published in 1989 and
a lot of resources easily available now were not within a student's reach four
decades ago. I hope to make plenty of digressions as I study these chapters.

\section{Notation}
Dynamical variables are always operators in quantum mechanics. Therefore, we will
not use a separate notation for operators. We will also use the same symbol for
its matrix representation. Whether a symbol means the abstract operator or its
matrix representation is clear from the context.

\section{The variation method}
We want to find an approximate solution to the time-independent Schr\"{o}dinger
equation
\begin{equation}\label{c1e1}
H\ket{\Phi} = E\ket{\Phi}
\end{equation}
Suppose that $H$ has a discrete spectrum of eigenvalues $\{E_n\}$
with $n = 0, 1, \ldots$. Let $\{\ket{\Phi_n}\}$ be the corresponding eigenkets.
Since $H$ is hermitian, the eigenvalues are real and the eigenkets can be
assumed to be orthonormal. Therefore, we also have,
\begin{equation}\label{c1e2}
\braket{\Phi_a|H|\Phi_b} = E_a\delta_{ab}.
\end{equation}
We can also assume that $\{\ket{\Phi_a}\}$ form a complete set so that any other
ket $\ket{\tilde{\Phi}}$ obeying the same boundary conditions can be expressed
as 
\begin{equation}\label{c1e3}
\ket{\tilde{\Phi}} = \sum_a\ket{\Phi_a}\braket{\Phi_a|\tilde{\Phi}}.
\end{equation}

In almost all cases, except the simplest system, $\{\ket{\Phi_a}\}$ are unknown.
But their existence can be assumed following the mathematical properties of
$H$. Therefore, \ref{c1e3} is true in principle. If $\ket{\tilde{\Phi}}$
is an approximate solution to \eqref{c1e1} then we will show that
\begin{equation}\label{c1e4}
\braket{\tilde{\Phi}|H|\tilde{\Phi}} \ge E_0.
\end{equation}
That is, the true ground state energy of the system is a lower bound to the
expectation value of $H$ in the state $\ket{\tilde{\Phi}}$. Therefore, in
any method to find an approximate solution to \eqref{c1e1} we choose the ket
that minimises the expectation value of $H$ in it. Typically, we consider
a family of \emph{trial kets} that depend on a set of parameters. We choose the
parameters that minimise the expectation value of the hamiltonian in the state.
Before illustrating this method for a few simple problems, we will prove the
inequality \eqref{c1e4}. Let
\begin{equation}\label{c1e5}
\ket{\tilde{\Phi}} = \sum_a c_a\ket{\Phi_a}
\end{equation}
so that,
\[
H\ket{\tilde{\Phi}} = \sum_a c_aE_a\ket{\Phi_a}
\]
and
\begin{equation}\label{c1e6}
\braket{\tilde{\Phi}|H|\tilde{\Phi}} = \sum_{a,b}c_b^\ast c_aE_a
\braket{\Phi_b|\Phi_a} = \sum_a |c_a|^2E_a.
\end{equation}
Since $\ket{\tilde{\Phi}}$ can be assumed to be normalised. 
\[
\sum_a |c_a^2| = 1.
\]
If the eigenstates are labelled such that $E_0 \le E_1 \le
\ldots$ then
\[
\braket{\tilde{\Phi}|H|\tilde{\Phi}} \ge E_0\sum_a|c_a|^2 =
E_0.
\]

Examples:
\begin{enumerate}
\item This is exercise 1.18 in the book \cite{so1989}. We will first solve
the hamiltonian analytically. Let
\[
\left(-\frac{1}{2}\frac{d^2}{dx^2} - \delta(x)\right)\psi(x)=E\psi(x).
\]
In the region $x \ne 0$ the solution is
\[
\psi(x)=\begin{dcases}Ae^{-\kappa x} & x > 0 \\Be^{\kappa x} & x < 0\end{dcases}
\]
where $A$ and $B$ are constants. Since $\psi$ must be continuous at $x = 0$,
we must have $A = B$. This gives us the energy eigenvalue
\[
E = -\frac{\kappa^2}{2}.
\]
If we integrate the Schr\"{o}dinger equation over the region $[-a, a]$ then
we have
\[
-\frac{1}{2}(\psi^\op(a) - \psi^\op(-a)) - \psi(0) = E
\int_{-a}^a\psi(x)dx.
\]
so that
\[
-\frac{1}{2}(-\kappa A e^{-\kappa a} - \kappa A e^{+\kappa a}) = 
\psi(0) + E\int_{-a}^a\psi(x)dx.
\]
In the limit $a \rightarrow 0$, since $\psi$ is continuous, the integral 
vanishes and this equation gives
\[
\kappa = 1
\]
Thus the true energy eigenvalue is $-1/2$.

We will now use the variational method to get it. Let $\tilde{\Phi} = Ne^{-ax^2}$
be a trial function. The normalisation condition gives 
\[
N^2 = \sqrt{\frac{2a}{\pi}}.
\]
We next evaluate
\[
H\tilde{\Phi} = aN(1 - 2ax^2)e^{-ax^2} - \delta(x)Ne^{-ax^2}.
\]
so that
\[
\braket{H}(a) = \int_{-\infty}^\infty dx\left(aN^2(1 - 2ax^2)e^{-2ax^2} - 
\delta(x)N^2 e^{-2ax^2}\right) = \frac{a}{2} - \sqrt{\frac{2a}{\pi}}.
\]
This function has an extremum at $a_0 = 2/\pi$ so that
\[
\braket{H}(a_0) = -\frac{1}{\pi} \approx -0.3183 > -0.5.
\]
We used equation \eqref{a1e8} to evaluate the integrals in this example.

\item This is exercise 1.19 in the book \cite{so1989}. The hamiltonian for the
hydrogen atom is
\[
H = \left(-\frac{1}{2}\nabla^2 - \frac{1}{r}\right).
\]
Once again we try the gaussian trial function
\[
\ket{\tilde{\Phi}} = Ne^{-ar^2}.
\]
This is a three dimensional wave-function and therefore the normalisation
condition is
\[
\int_0^\infty\int_0^{2\pi}\int_0^\pi N^2 e^{-2ar^2}r^2\sin\theta drd\theta d\phi
= 4\pi N^2 \frac{1}{8a}\sqrt{\frac{\pi}{2a}} = 1
\]
so that
\[
N^2 = \sqrt{\frac{2a}{\pi}}
\]
Now,
\[
H\ket{\tilde{\Phi}} = aN(3 - 2ar^2)e^{-ar^2} - N\frac{e^{-ar^2}}{r}
\]
so that
\[
\braket{H}(a) = \frac{3a}{2} - 2\sqrt{\frac{2a}{\pi}}.
\]
Its minimum is as 
\[
a_0 = \frac{8}{9\pi}
\]
so that the minimum expectation value is
\[
\braket{H}(a) = -\frac{4}{3\pi} \approx -0.4244,
\]
while the actual ground state energy is $-0.5$ hartree.
\end{enumerate}

In the two examples considered above the trial function was a gaussian. In most
hamiltonians of interest, the trial function must have several parameters and 
they are optimised to minimise the expectation value. This procedure is 
simplified if we consider a set of ``basis" functions $\{\ket{\Psi_i}\}$, each 
of which is individually well-understood and construct a trial function of the 
form
\begin{equation}\label{c1e7}
\ket{\tilde{\Phi}} = \sum_{i=1}^N c_i\ket{\Psi_i}.
\end{equation}
We assume, for now, that the basis functions are real-valued and orthonormal.
Let
\begin{equation}\label{c1e8}
H_{ij} = \braket{\Psi_i|H|\Psi_j}.
\end{equation}
$H$ being hermitian and $\{\ket{\Psi_i}\}$ being real-valuef guarantees
that
\begin{equation}\label{c1e9}
H_{ij} = H_{ji}.
\end{equation}
We also assumet that $\ket{\tilde{\Phi}}$ is normalised so that
\begin{equation}\label{c1e10}
\sum_{i=1}^N |c_i|^2 = 1.
\end{equation}
The variational procedure requires us to minimise
\begin{equation}\label{c1e11}
\braket{\tilde{\Phi}|H|\tilde{\Phi}} = \sum_{i,j}c_ic_jH_{ij},
\end{equation}
by treating it as a function if $N$ parameters $c_1, \ldots, c_N$. They cannot
be chosen arbitrarily because they must satisfy \eqref{c1e10}. We use Lagrange's
method of undetermined multiplier to solve this minimisation problem. If
\begin{equation}\label{c1e12}
\mathscr{L}(c_1, \ldots, c_N, E) = \braket{\tilde{\Phi}|H|\tilde{\Phi}} 
- E(\braket{\tilde{\Phi}|\tilde{\Phi}} - 1),
\end{equation}
where $E$ is the undetermined multiplier then we find the minimum using the
conditions
\begin{equation}\label{c1e13}
\pdt{\mathscr{L}}{c_k} = 0, k = 1, \ldots, N - 1.
\end{equation}
We then choose $E$ such that
\begin{equation}\label{c1e14}
\pdt{\mathscr{L}}{c_N} = 0.
\end{equation}
Using \eqref{c1e10} and \eqref{c1e11} in \eqref{c1e12}, we get
\begin{equation}\label{c1e15}
\mathscr{L}(c_1, \ldots, c_N, E) = \sum_{i,j}c_ic_jH_{ij} - E\left(\sum_i|c_i|^2
- 1\right),
\end{equation}
so that
\[
\pdt{\mathscr{L}}{c_k} = \sum_jc_jH_{kj} + \sum_ic_iH_{ik} - 2Ec_k.
\]
Using \eqref{c1e9}, it can be simplified to
\begin{equation}\label{c1e16}
\pdt{\mathscr{L}}{c_k} = 2\sum_i H_{kj}c_j - 2Ec_k.
\end{equation}
If $\vec{c} = (c_1, \ldots, c_N)^T$ then we can express the minimisation 
condition as the matrix eigenvalue equation
\begin{equation}\label{c1e17}
H\vec{c} = E\vec{c}.
\end{equation}
Since $H$ is symmetric and real-valued, the eigenvalues $E_1, \ldots, E_N$ are
real. We label them so that $E_1 \le E_2 \le \ldots \le E_N$. Let $\vec{c}^1,
\ldots \vec{c}^N$ be the corresponding eigenvectors which can be chosen to be
orthonormal. Each of these eigenvectors gives the $N$ coefficients in 
\eqref{c1e7}. Thus, we get $N$ trial functions
\begin{equation}\label{c1e18}
\ket{\Phi^a} = \sum_i c^a_i\ket{\Psi_i}.
\end{equation}
They are guaranteed to be orthonormal because
\begin{equation}\label{c1e19}
\braket{\Phi^a|\Phi^b} = \sum_{ij}c^a_ic^b_j\braket{\Psi_i|\Psi_j} = 
\sum_i c^a_ic^b_j = \delta_{ab},
\end{equation}
where in the second term we used orthonormality of the basis functions and in
the third term we used the orthonormality of $\{\vec{c}^i\}$. The equation
\begin{equation}\label{c1e20}
\braket{\tilde{\Phi}^a|H|\tilde{\Phi}^a} = \sum_{i,j}c_i^ac_j^a\braket{
\Psi_i|H|\Psi_j} = \sum_{i,j}c_i^a H_{ij} c_j^a = E_a\sum_i{c_i^a}^2
= E_a.
\end{equation}
Thus, the eigenvalues of \eqref{c1e17} are the expectation values of the 
hamiltonian in the $N$ trial states.

We shall demonstrate this procedure to find the energy levels of a hydrogen
atom in an external electric field $F\uv{z}$. The hamiltonian, in atomic units,
is
\begin{equation}\label{c1e21}
H = -\frac{1}{2}\nabla^2 - \frac{1}{r} + Fr\cos\theta = H_0 +
Fr\cos\theta.
\end{equation}
We choose the basis functions $\{\ket{1s}, \ket{2p_z}\}$ so that the trial
function is
\begin{equation}\label{c1e22}
\ket{\tilde{\Phi}} = c_1\ket{1s} + c_2\ket{2p_z}.
\end{equation}
We also have
\begin{eqnarray*}
H_0\ket{1s} &=& -\frac{1}{2}\ket{1s} \\
H_0\ket{2p_z} &=& -\frac{1}{8}\ket{2p_z}
\end{eqnarray*}
Our first task is to compute the matrix elements $H_{ij}$.
\begin{equation}\label{c1e23}
H_{11} = \braket{1s|H_0 + Fr\cos\theta|1s} = -\frac{1}{2} + F\int_0^\infty
\int_0^\pi\int_0^{2\pi} \frac{e^{-2r}}{\pi}r^3\cos\theta\sin\theta drd\theta
d\phi,
\end{equation}
where we used
\[
\ket{1s} = \frac{e^{-r}}{\sqrt{\pi}}.
\]
The $\theta$-integral can be readily confirmed to vanish so that
\begin{equation}\label{c1e24}
H_{11} = -\frac{1}{2}.
\end{equation}

We now compute
\begin{equation}\label{c1e25}
H_{22} = -\frac{1}{8} + \frac{F}{32\pi}\int_0^\infty\int_0^\pi\int_0^{2\pi}
r^4e^{-r}\cos^3\theta\sin\theta drd\theta d\phi,
\end{equation}
where we used
\begin{equation*}
\ket{2p_z} = \frac{1}{\sqrt{32\pi}} re^{-r/2}\cos\theta.
\end{equation*}
Once again the $\theta$ integral vanishes and we get
\begin{equation}\label{c1e26}
H_{22} = -\frac{1}{8}F.
\end{equation}

The off-diagonal term is
\begin{equation}\label{c1e27}
H_{12} = -\frac{1}{8}\braket{1s|2p_z} + \frac{F}{\pi\sqrt{32}}\int_0^\infty
\int_0^\pi\int_0^{2\pi} r^3 e^{-3r/2}\cos^2\theta\sin\theta drd\theta d\phi
\end{equation}
The first term vanishes because of the orthogonality of the $1s$ and $2p_z$ 
states. After evaluating the integral in the second term,
\begin{equation}\label{c1e28}
H_{12} = F\frac{128\sqrt{2}}{243}
\end{equation}

Let us introduce the constant
\begin{equation}\label{c1e29}
\beta = \frac{128\sqrt{2}}{243}
\end{equation}
so that the $H$-matrix is
\begin{equation}\label{c1e30}
H = \begin{bmatrix} -1/2 & F\beta \\ F\beta & -1/8 \end{bmatrix}.
\end{equation}
Its secular equation is
\begin{equation}\label{c1e31}
E^2 +  (H_{11} + H_{22})E + (H_{11}H_{22} - H_{12})^2 = 0
\end{equation}
so that the eigen values are
\begin{equation}\label{c1e32}
E = \frac{H_{11} + H_{22}}{2} \pm \sqrt{\left(\frac{H_{11} - H_{22}}{2}\right)^2
+ H_{12}^2}.
\end{equation}
Substituting the values of the matrix elements, we get
\begin{equation}\label{c1e33}
E = -\frac{5}{16} \pm \sqrt{\frac{9}{256} + \frac{32768}{59049}F^2}.
\end{equation}
When the field is weak, that is $F$ is small, we can approximate this expression
as
\[
E = -\frac{5}{16} \pm \frac{3}{16}\left(1 + \frac{1}{2}\frac{256}{9}
\frac{32768}{59049}F^2\right)
\]
so that
\[
E_- = -\frac{1}{2} - \frac{1}{2}\frac{3}{16}\frac{256}{9}\frac{32768}{59049}F^2
= -\frac{1}{1} - \frac{1}{2}(2.9596 F^2).
\]

\section{Standard integrals}
The most basic Gaussian integral is
\begin{equation}\label{a1e1}
\int_{-\infty}^\infty e^{-x^2}dx = \sqrt{\pi}.
\end{equation}
To evaluate it, start with 
\[
I^2 = \int_{-\infty}^\infty\int_{-\infty}^\infty e^{-x^2-y^2}dxdy
\]
Transforming to polar coordinates
\[
I^2 = \int_0^\infty\int_0^{2\pi} e^{-r^2}rdrd\theta = 
2\pi\int_0^\infty e^{-r^2}rdr = \pi
\]
from which equation \eqref{a1e1} follows immediately. From \eqref{a1e1} we 
have, using a quick transformation of variables
\begin{equation}\label{a1e2}
\int_{-\infty}^\infty e^{-ax^2}dx = \sqrt{\frac{\pi}{a}}.
\end{equation}
The integral
\[
\int_{-\infty}^\infty e^{-ax^2+bx}dx
\]
is evaluated by completing the square
\[
-ax^2+bx = -\left(ax^2 - 2\sqrt{a}x\frac{b}{2\sqrt{a}} + \frac{b^2}{4a}\right)
+ \frac{b^2}{4a} = -\left(ax - \frac{b}{2\sqrt{a}}\right)^2 + \frac{b^2}{4a}
\]
we have
\[
\int_{-\infty}^\infty e^{-ax^2+bx}dx = \exp\left(\frac{b^2}{4a}\right)
\int_{-\infty}^\infty\exp\left[-\left(ax-\frac{b}{2\sqrt{a}}\right)^2\right]dx
\]
so that
\begin{equation}\label{a1e3}
\int_{-\infty}^\infty e^{-ax^2+bx}dx = 
\exp\left(\frac{b^2}{4a}\right)\sqrt{\frac{\pi}{a}}.
\end{equation}
A trivial generalisation of this result is
\begin{equation}\label{a1e4}
\int_{-\infty}^\infty e^{-ax^2+bx+c}dx = \exp\left(\frac{b^2}{4a} + c\right)
\sqrt{\frac{\pi}{a}}.
\end{equation}
The integral \eqref{a1e1} is related to the gamma function. For, using the 
transformation $t = x^2$,
\[
\int_{-\infty}^\infty e^{-x^2}dx = 2\int_{0}^\infty e^{-x^2}dx =
\int_{0}^\infty e^{-t}t^{-1/2}dt = \Gamma\left(\frac{1}{2}\right)
\]
where
\begin{equation}\label{a1e5}
\Gamma(z) = \int_{0}^\infty t^{z-1}e^{-t}dt.
\end{equation}
Thus we have the frequently used result
\begin{equation}\label{a1e6}
\Gamma\left(\frac{1}{2}\right) = \sqrt{\pi}.
\end{equation}
Integrating \eqref{a1e5} by parts once we get
\begin{equation}\label{a1e7}
\Gamma(z) = (z - 1)\Gamma(z - 1).
\end{equation}

We now start with equation \eqref{a1e2} and differentiate it with respect to 
$a$ to get
\[
\int_{-\infty}^\infty -x^2e^{-ax^2}dx = -\frac{1}{2}\frac{\sqrt{\pi}}{a^{3/2}}
\]
or
\[
\int_{-\infty}^\infty x^2e^{-ax^2}dx = 
\frac{1}{2}\Gamma\left(\frac{1}{2}\right)\frac{1}{a^{3/2}} = 
\frac{\Gamma(3/2)}{a^{3/2}}.
\]
Integrating it once more with respect to $a$ you get
\[
\int_{-\infty}^\infty x^4e^{-ax^2}dx = \frac{\Gamma(5/2)}{a^{5/2}}.
\]
Integrating it $n$ times, we get the result
\begin{equation}\label{a1e8}
\int_{-\infty}^\infty x^{2n}e^{-ax^2}dx = 
\frac{\Gamma(\frac{2n+1}{2})}{a^{(2n+1)/2}}.
\end{equation}
and hence
\begin{equation}\label{a1e9}
\int_{0}^\infty x^{2n}e^{-ax^2}dx = 
\frac{1}{2}\frac{\Gamma(\frac{2n+1}{2})}{a^{(2n+1)/2}}.
\end{equation}
It is clear that
\begin{equation}\label{a1e10}
\int_{-\infty}^\infty x^{2n+1}e^{-ax^2}dx = 0.
\end{equation}
However, we can the integral with limits $0$ to $\infty$ will be non-zero. To 
get their value, we start with
\[
\int_0^\infty x e^{-ax^2}dx = \frac{1}{2a}
\]
and differentiate it $n$ times to get
\begin{equation}\label{a1e11}
\int_0^\infty x^{2n+1} e^{-ax^2}dx = \frac{n!}{2a^{n+1}}.
\end{equation}
By putting $2n+1=k$ in equation \eqref{a1e9} and $2n=k$ in equation 
\eqref{a1e11} we get the single formula
\begin{equation}\label{a1e12}
\int_0^\infty x^{k} e^{-ax^2}dx = 
\frac{\Gamma\left(\frac{k+1}{2}\right)}{2a^{(k+1)/2}}.
\end{equation}

\section{Constrained optimisation}
Consider the problem of finding the maximum of a function $f:\mathbb{R}^2
\rightarrow \mathbb{R}$ subject to the constraint $g(\vec{r}) = c$ where
$g:\mathbb{R}^3 \rightarrow \mathbb{R}$. If there was no constraint, we
would have equated $\grad f = 0$. However the constraint requires us to traverse
the surface $z = f(x, y)$ along the curve $g(x, y, z) = c$. In general, 
travelling along the curve $g(x, y, z) = c$ involves crossing the level curves
of $f$. However, at the extrema, the level curves of $f$ are tangential to $g$
or that the gradients of the two functions are parallel to each other. Therefore,
there exists a constant $\lambda$ such that $\grad f = \lambda\grad g$. 
Equivalently, the condition for constrained optimisation is
\begin{equation}\label{a2e1}
\grad (f - \lambda g) = 0.
\end{equation}
The constant $\lambda$ need not be calculated and is therefore called an
undetermined multiplier of the constraint.
